\section{Time stepping and diagnostics in slice variables}\label{app:timestepping}

This appendix analyses GANDALF's two time-stepping schemes (version 0.6.0), the implicit--explicit scheme IMEX-RK222, which is the scheme ARS(2,2,2) of \citet{AscherRuuthSpiteri1997}, and the integrating-factor scheme Lawson-RK4 \citep{Lawson1967}, in the slice variables of theorem~\ref{thm:trunc}, and lists the factors that enter its Hermite diagnostics. We write $\Delta t$ for the time step, $k_{\perp,\max}$ and $k_{z,\max}$ for the largest perpendicular and parallel wavenumbers of the dealiasing band, and
\begin{equation*}
  N(g)=-\{\Phi,g\}-s\{\Psi,Ag\}
\end{equation*}
for the part of the right-hand side of (\ref{tr:hierarchy}) that both schemes treat explicitly. By theorem~\ref{thm:trunc} its projection on slice $j$ is $N_j(a_j)=-\{\Phi+su_j\Psi,a_j\}$. All tests through GANDALF's time step quoted below are short integrations in double precision.

\subsection{The time step in slice variables}\label{appD:step}

For each Fourier mode let $\mathsf L=-\ii sk_zA-\nu\mathcal D$, where $\mathcal D=\mathrm{diag}\big((m/M)^n\big)$ for $m\ge2$ and $\mathcal D_{00}=\mathcal D_{11}=0$ describe the hypercollisions. GANDALF's IMEX-RK222 step uses $\gamma=1-1/\sqrt2$ and $\delta=1-1/(2\gamma)=-1/\sqrt2$:
\begin{align*}
  (I-\gamma\Delta t\,\mathsf L)\,g^{(1)}&=g^n+\gamma\Delta t\,N(g^n;z^\pm_n),\\
  (I-\gamma\Delta t\,\mathsf L)\,g^{*}&=g^n+(1-\gamma)\Delta t\,\mathsf L\,g^{(1)}+\delta\Delta t\,N(g^n;z^\pm_n)+(1-\delta)\Delta t\,N\big(g^{(1)};z^\pm_{1/2}\big),
\end{align*}
followed by $g^{n+1}=\mathcal P\mathcal K g^*$, where $\mathcal P$ is the dealiasing mask, the projection on the dealiasing band, and
\begin{equation*}
  \mathcal K=\exp\big[-\eta\,(k_\perp^2/k_{\perp,\max}^2)^r\,\Delta t\big]\,\exp\big[-\eta_z\,(k_z^2/k_{z,\max}^2)^{r_z}\,\Delta t\big];
\end{equation*}
both act on every moment in the same way. Here $z^\pm_n$ are the Elsasser potentials at $t_n$, and $z^\pm_{1/2}=e^{\pm\ii k_z\Delta t/2}[z^\pm_n+\frac12\Delta t\,\mathcal N^\pm_n]$ is GANDALF's prediction of them at $t_n+\Delta t/2$, with $\mathcal N^\pm_n$ the nonlinear terms of the Elsasser equations at $t_n$. Without dissipation and hypercollisions, a step reassembled from krmhd's own functions in this order reproduces GANDALF's step exactly, with zero difference.

In slice variables and with $\nu=0$, $\mathsf L$ acts on slice $j$ as multiplication by $-\ii sk_zu_j$, $N$ acts as $N_j$, and $\mathcal P\mathcal K$ acts on every slice as it acts on every moment. Each slice is therefore advanced by the same scheme applied to the scalar problem $\partial_ta_j=-\ii sk_zu_ja_j+N_j(a_j)$, with no reference to the other slices. With $\nu>0$, $\mathsf L$ becomes $-\ii sk_z\,\mathrm{diag}(u_j)-\mathsf C$ in the slice basis, with $\mathsf C$ from (\ref{tr:coll}), and the implicit solves couple the slices. The Lawson-RK4 step applies the exact streaming factor $\exp(-\ii sk_zA\Delta t)$, computed in the eigenbasis of $A$, takes four RK4 stages of $N$ with $z^\pm$ at $t_n$, $t_n+\Delta t/2$ (twice) and $t_n+\Delta t$, and then applies the factor $\exp(-\nu\mathcal D\Delta t)$, the two damping factors of $\mathcal K$ and the mask. In slice variables the streaming factor is $e^{-\ii sk_zu_j\Delta t}$, and again only the hypercollisions couple the slices. GANDALF's routine \texttt{compute\_streaming\_\allowbreak eigensystem} returns left eigenvectors parallel to $\boldsymbol w_j=\Pw\boldsymbol r_j$, to $4\times10^{-16}$ at $M=64$ and 65, so the slice amplitudes can be computed from GANDALF's own eigensystem.

\subsection{Stability of the explicit stages}\label{appD:stab}

For the test equation $\dot y=\Omega y$, with $\zeta=\Omega\Delta t$, the explicit stages of ARS(2,2,2) give $y^{(1)}=(1+\gamma\zeta)\,y$ and $y^{n+1}=[1+\delta\zeta+(1-\delta)\zeta(1+\gamma\zeta)]\,y$. Since $\gamma(1-\delta)=\frac12$, the stability function is
\begin{equation}\label{appD:R2}
  R_2(\zeta)=1+\zeta+\frac{\zeta^2}{2},\qquad |R_2(\ii y)|^2=1+\frac{y^4}{4}.
\end{equation}
The four stages of Lawson-RK4 at $k_z=0$ are those of classical RK4, with $|R_4(\ii y)|^2=1-y^6/72+y^8/576$, which does not exceed 1 if and only if $|y|\le2\sqrt2$. We derived both with GANDALF's constants in SymPy. The implicit part does not remove the instability of (\ref{appD:R2}): in the frozen-line test described below, including the implicit streaming in the stability function of the full step raises the largest amplification per step from $1.92$ to $2.16$.

Let $\Phi$ and $\Psi$ lie inside the dealiasing band and be independent of time. On the band, $N_j$ is then antisymmetric in the $L^2$ product of real fields, because the pseudo-spectral bracket of band-limited fields is antisymmetric under the volume average. It is therefore a normal operator with imaginary eigenvalues $\pm\ii\omega$, and its spectral radius $\rho_j=\max|\omega|$ equals its norm. The pseudo-spectral bracket of band-limited fields is the exact bracket projected on the band, so $\|N_ja\|\le V_j\|\nabla_\perp a\|\le k_{\perp,\max}V_j\|a\|$, with $V_j=\max|\nabla_\perp(\Phi+su_j\Psi)|$, and $\rho_j\le k_{\perp,\max}V_j$. For a normal operator, $\|p(\Delta t\,N_j)\|$ is the largest value of $|p(\ii\omega\Delta t)|$ over its eigenvalues, for any polynomial $p$. When $\Phi$ and $\Psi$ do not depend on $z$, $N_j$ commutes with the streaming factor of slice $j$, and the following statements are exact. Without dissipation and hypercollisions, the Lawson-RK4 step does not amplify slice $j$ if and only if $\Delta t\,\rho_j\le2\sqrt2$; damping that commutes with $N_j$ can only enlarge this range. On the plane $k_z=0$, where the implicit part acts only through dissipation, the IMEX step without dissipation amplifies the fastest-rotating mode of slice $j$ by $[1+(\Delta t\rho_j)^4/4]^{1/2}$ per step, at every $\Delta t$. If dissipation multiplies the slice by $e^{-d_j\Delta t}$ per step and commutes with $N_j$, the step is stable there if and only if $1+(\Delta t\rho_j)^4/4\le e^{2d_j\Delta t}$, which for small $\Delta t\rho_j$ is (\ref{eq:explicitstability}). Hypercollisions add approximately their rate $\mathsf C_{jj}$ to $d_j$, and for the fastest slices this rate is between $0.63\nu$ and $0.90\nu$ for $n=6$ and $2$ (section~\ref{tr:collsec}). When the fields depend on $z$ or on time, or the dissipation does not commute with $N_j$, these conditions are estimates.

GANDALF's hyper-resistive factor does not commute with $N_j$ unless the flow is uniform. To estimate its effect we use a local model, in which a Fourier mode of wavenumber $k_\perp$ parallel to a locally uniform flow of speed $V_j$ rotates at $\omega=k_\perp V_j$ and is damped at the rate $d(k_\perp)=\eta\,(k_\perp/k_{\perp,\max})^{2r}$. For small $y=k_\perp V_j\Delta t$, (\ref{eq:explicitstability}) with $\rho_j=k_\perp V_j$ and $d_j=d(k_\perp)$ becomes
\begin{equation*}
  \Delta t^3\le\frac{8\eta}{V_j^4}\,\frac{k_\perp^{2r-4}}{k_{\perp,\max}^{2r}} .
\end{equation*}
For $r=2$ the right-hand side does not depend on $k_\perp$, and for $r=1$ it is smallest at $k_\perp=k_{\perp,\max}$. In both cases
\begin{equation}\label{appD:thr}
  \Delta t\le\Delta t_{\mathrm c}(V_j)=\frac{(8\eta)^{1/3}}{(k_{\perp,\max}V_j)^{4/3}}\qquad(r\le2).
\end{equation}
For $r>2$ the smallest perpendicular wavenumber $k_{\perp,\min}$ of the grid binds, and the threshold is smaller than (\ref{appD:thr}) by the factor $(k_{\perp,\max}/k_{\perp,\min})^{(2r-4)/3}$, which is $(k_{\perp,\max}/k_{\perp,\min})^{4/3}$ for $r=4$. In a computation of the per-mode amplification of this model on a $128^2$ grid, the exact threshold agrees with (\ref{appD:thr}) to $0.2\%$ for $r=1$ and $r=2$, and for $r=4$ the smallest exact threshold is reached at the smallest wavenumber of the computation and agrees with the reduced value for that wavenumber.

We tested these statements through GANDALF's time step in three settings.

The first setting is a stationary shear flow, $\Phi=0.3a\cos2\pi y$ and $\Psi=a\cos2\pi y$, with $a$ chosen so that $\max|\nabla_\perp\Psi|=2$, on a $32\times32\times2$ grid with $M=64$, $\Lambda=\sqrt5$ and $s=1$. GANDALF's Elsasser step is exact for these fields, and each slice sees the fixed operator $N_j=\ii k_xf_j'(y)$, $f_j=\Phi+su_j\Psi$, whose spectral radius on the plane $k_x=k_{x,\max}$ is at most $k_{x,\max}V_j$. The trajectory of every slice through GANDALF's step equals the predicted per-slice linear map to $7\times10^{-14}$, for IMEX with $\eta=0$ and $\eta=2000$ and for Lawson-RK4 below and above its limit. The fastest slice is advected at $V_j=21.8$, while the speed in GANDALF's time-step estimate is $2.6$, and the estimate's $\Delta t=3.61\times10^{-3}$ corresponds to $y=k_{x,\max}V_j\Delta t=4.9$ for that slice. Without dissipation the slice that grows fastest is the one with $|u_j|=u_{\max}$. Lawson-RK4 is stable at $y=2.6$ for the fastest slice and unstable at $y=3.1$, where that slice grows by a factor $6\times10^4$ in 20 steps. With $\eta=2000$ and $r=2$, the IMEX scheme becomes unstable at $\Delta t=1.25\times10^{-3}$ for the slice at $u_{\max}$, where (\ref{appD:thr}) gives $1.06\times10^{-3}$, and at $\Delta t=0.21$ for the slice at $u=0$. At GANDALF's $\Delta t$, 15 steps multiply the slice at $u_{\max}$ by 765, while the slice at $u=0$ decays by the factor $3\times10^{-13}$.

The second setting is a random three-dimensional state on a $16^3$ grid, with $M=64$, $\Lambda=\sqrt5$, $\beta_i=1$ and random $z^\pm$ with $\max|\nabla_\perp z^+|=5$ and $\max|\nabla_\perp z^-|=4.4$, so that the speed in GANDALF's estimate is set by $|\nabla_\perp z^\pm|$. The fastest slice is advected at $V_j=32.1$, and $k_{\perp,\max}V_j\Delta t=5.35$ at GANDALF's $\Delta t$. The spectral radius of its advection operator, computed by power iteration, is $\rho_j=414$, which is $0.29\,k_{\perp,\max}V_j$, so that $\Delta t\,\rho_j=1.55$ at GANDALF's $\Delta t$. Lawson-RK4 is stable there. It is also stable at $\Delta t\,\rho_j=2.40$ and unstable at $3.54$, where the energy of the fast slices grows by a factor $7\times10^{29}$ in 40 steps, as the limit $\Delta t\,\rho_j\le2\sqrt2$ predicts. IMEX at GANDALF's $\Delta t$, without dissipation, multiplies the energy of the slices with $|u_j|>8$ by $4.6\times10^{12}$ in 30 steps, while that of the slices with $|u_j|<2$ changes by less than 1\%. Whether GANDALF's estimate is stable for Lawson-RK4 therefore depends on the state, through the ratio $\rho_j/(k_{\perp,\max}V_j)$. For the IMEX scheme without dissipation no time step makes the explicit part stable, and the state sets only the growth rate.

The third setting is the frozen-line test. We advanced the initial state of the test of corollary~\ref{cor:singlepoint} with GANDALF's time step: random fields with $\mathrm{rms}\,|\nabla_\perp\Phi|=\mathrm{rms}\,|\nabla_\perp\Psi|=1$ on a $32^3$ grid, $M=8$, $\Lambda=\sqrt5$ and $s=1$, with no dissipation or collisions. The nominal values of $k_{\perp,\max}V_j\Delta t$ for the fastest slice are $1.81$ at $\Delta t=0.002$ and $0.91$ at $\Delta t=0.001$, for which (\ref{appD:R2}) gives $|R_2|=1.92$ and $1.08$ per step. These values use the full advecting potential $\Phi+su_j\Psi$; the potential $su_j\Psi$ alone gives $1.51$ and $0.75$. At $\Delta t=0.002$ the IMEX run is in exponential growth by $t=0.128$, where the free energy $\sum_mp_m\langle g_m^2\rangle$, which the equations conserve, is $2.5$ times its initial value. It is 96 times its initial value at $t=0.144$ and $9.3\times10^6$ times at $t=0.208$. At $t=0.128$ the largest Fourier coefficient of $g$ is the same as in the stable runs. The growth shows in the conserved free energy and near the cutoff, where the share of the energy of $g$ in the outermost shell of wavenumbers rises from $2.8\times10^{-2}$ at $t=0.112$ to $0.99$ at $t=0.144$. Lawson-RK4 at the same $\Delta t$, and IMEX at $\Delta t=0.001$, conserve the free energy up to $t=0.4$. At $\Delta t=0.001$ no growth is visible, although the nominal factor $1.08$ per step would amplify the fastest mode by many orders of magnitude in 400 steps, so the nominal value overstates the actual growth rate. In the stable runs the IMEX step converges in $\Delta t$ at order $1.01$ in $g$ and $2.00$ in $z^+$ at $t=0.064$, and Lawson-RK4 at order $2.00$ in $g$.

In the first two settings the slices that limit the time step are the fastest ones. A speed that bounds $V_j$ for every slice is $\max|\nabla_\perp\Phi|+su_{\max}\max|\nabla_\perp\Psi|$, with $u_{\max}$ from table~\ref{tr:umaxtab}. A corrected speed in the time-step estimate is not sufficient for the IMEX scheme, however. Its explicit part has no stable interval on the imaginary axis, so its stability rests on dissipation through (\ref{eq:explicitstability}). The alternatives are a smaller $\Delta t$ or an explicit tableau whose stability region contains an interval of the imaginary axis.

\subsection{The coupling of \texorpdfstring{$g$}{g} to the Elsasser fields}\label{appD:order}

The explicit tableau of ARS(2,2,2) has the weights $\tilde b=(\delta,1-\delta,0)$ and the stage times $\tilde c=(0,\gamma,1)$, so that $\sum\tilde b=1$ and $\sum\tilde b\tilde c=(1-\delta)\gamma=\frac12$, which are the conditions for second order. The explicit term $N$ depends on time through $z^\pm(t)$. GANDALF evaluates the second stage with $z^\pm$ predicted to $t_n+\Delta t/2$ rather than to $t_n+\gamma\Delta t$, so for this dependence the effective condition is
\begin{equation*}
  (1-\delta)\cdot\frac12=\frac{1}{4\gamma}=\frac12+\frac{\sqrt2}{4}\neq\frac12 .
\end{equation*}
The local error then contains the term $\frac{\sqrt2}{4}\Delta t^2\,(\partial N/\partial z^\pm)\,\dot z^\pm$, and the global error of $g$ is first order in $\Delta t$. For the scalar model $\dot g=\ii\omega(t)\,g$ with $\omega=\omega_0+\omega_1t$, the step with stage time $\frac12$ has the local error $\ii\frac{\sqrt2}{4}\omega_1\Delta t^2\,g$, a pure phase error, while the step with stage time $\gamma$ has no error at order $\Delta t^2$. Both change $|g|^2$ only at order $\Delta t^3$ per step, by the factors $1+\frac{\sqrt2}{2}\omega_0\omega_1\Delta t^3$ for stage time $\gamma$ and $1+\frac{1+\sqrt2}{2}\omega_0\omega_1\Delta t^3$ for stage time $\frac12$, up to higher orders. In slice variables the defect advects each slice with an incompressible velocity taken at the wrong time. It does not change the invariants at leading order, and for advection that depends on time the drift of the invariants is second order in $\Delta t$ for both stage times.

We tested this on an $8^3$ grid with $M=6$, $\Lambda=\sqrt5$ and $\max|\nabla_\perp z^\pm|=2$, up to $T=0.2$, against a Lawson-RK4 reference with 2048 steps. For $\Delta t=T/16$ to $T/256$, $g(T)$ from GANDALF's IMEX step converges at the observed orders $1.35$, $1.12$, $1.04$ and $1.01$ for a single $z^+$ wave with $z^-=0$, for which GANDALF's Elsasser step is exact, and at $1.18$, $1.06$, $1.02$ and $1.01$ for general $z^\pm$. Lawson-RK4 converges at order $4.0$ for the single wave and $2.0$ for general $z^\pm$, where the midpoint prediction of the Elsasser fields limits it. The reassembled IMEX step with $z^\pm$ predicted to $t_n+\gamma\Delta t$ instead converges at order $2.00$ in both cases. For the single wave the relative drift of the free energy converges at the local orders $2.55$, $2.38$, $2.23$ and $2.13$, for GANDALF's step and for the consistent variant alike. For general $z^\pm$ the drift changes sign between $\Delta t=T/32$ and $T/128$, and no order can be read from it. In a second test, on an $8^3$ grid with $M=8$, $\Lambda=-\sqrt5$, $\beta_i=2$ and a single $z^+$ wave, $g(T)$ converges at orders that fall from $1.92$ to $1.001$. The drift of the free energy has the local orders $3.3$, $3.5$ and $5.3$ before it changes sign near $\Delta t=T/128$, and approaches 2 after it ($1.89$ and $1.95$ for the two finest pairs of steps). A scan over coarse time steps can therefore suggest an order of the invariant drift that is higher than the true one. GANDALF's documentation states that the mismatch of the stage times is absorbed by the second-order truncation error of the scheme; these tests show that the error of $g$ is first order. Predicting $z^\pm$ to $t_n+\gamma\Delta t$ restores second order.

\subsection{Diagnostics}\label{appD:diag}

GANDALF's time-step estimate, \texttt{compute\_cfl\_timestep}, sets $\Delta t=0.3\min(\Delta x,\Delta y,\Delta z)/v_{\max}$, with $v_{\max}=\max(v_A,\max|\nabla_\perp z^\pm|,s\max|\nabla_\perp\Psi|)$ and $0.3$ the default safety factor; we reproduced it exactly. As shown above, it omits the factor $u_j$ of the slice speeds. The diagnostic \texttt{hermite\_flux} returns, for each Fourier mode, $-k_z[2(m+1)]^{1/2}\,\mathrm{Im}\,[g_{m+1}g_m^*]$, which is the $k_z$-based flux $\Pi^z_{m+1/2}(\boldsymbol k)$ of section~\ref{sec:decorrelation} divided by $s\,p_m$. It omits the field-line term $\{\Psi,\cdot\}$ of $\gradpar$, the factor $s$, and at $m=0$ the factor $p_0=1-1/\Lambda$. On GANDALF's right-hand side, $s$ times the diagnostic equals the $\partial_z$ part of $\Pi_{m+1/2}$ to $3\times10^{-16}$ for $m\ge1$, and at $m=0$ it differs from it by the factor $1/(1-1/\Lambda)$, which is $1.809$ at $\Lambda=\sqrt5$ and $0.691$ at $\Lambda=-\sqrt5$. Section~\ref{sec:decorrelation} gives the size of the omitted field-line part in strongly nonlinear states. The diagnostic \texttt{hermite\_free\_energy} uses the weight $1$ on $g_M^2$, so it is the conserved free energy for the zero closure but not for the copy closure (section~\ref{tr:formsec}). GANDALF computes in single precision unless 64-bit arithmetic is enabled in JAX, and its implicit operator is then also built in single precision. In that case GANDALF's computation satisfies the identities of section~\ref{sec:truncated} only to single-precision round-off. The omission of the slice speed from \texttt{compute\_cfl\_timestep}, the stage time of the IMEX step (section~\ref{appD:order}) and the omission of the field-line term from \texttt{hermite\_flux} have been reported in GANDALF's issue tracker.
